\section{Descent, immersed curvature and Morse heights}
\label{sec:parallel-investigations}

This section records the third program edition. Checked implications below specify their witnesses and their remaining geometric inputs. For
physical shell selection and the residual of a repair family, the targets
are Lemma~\ref{lem:master-quantiles} and
Corollary~\ref{cor:master-repair-residual}. A supplied shell is tested by
Proposition~\ref{prop:master-harmonic-annulus}. The following results use
the normalization and oriented smooth category of
Section~\ref{sec:master-variational}.

\begin{proposition}[A joint competitor on the original smooth type]
\label{prop:routes-joint-descent}
\lean{SmoothPoincare.jointCompetitorGain_le_excess, SmoothPoincare.excludesStrictJointCompetitor}
\leanok
\uses{lem:master-competitor}
Fix a nonempty set of admissible positive-class energies on \(M\) with
finite infimum \(I\). A source has energy \(e\). Suppose a marked
replacement of a region, including its transitions, supplies a smooth
positive metric on the original type, whose total energy is at most
\(e-m+\gamma\). Here \(m\) is the removed raw energy and \(\gamma\)
includes the whole replacement cost and any net change outside the
region. Then \(m-\gamma\le e-I\). A strict reverse inequality for one
actual competitor gives a contradiction.
\end{proposition}

\begin{proof}
\leanok
The marked pullback supplies membership on the original type. Subtract the
infimum bound from the complete energy comparison.
\end{proof}

\begin{lemma}[Weyl energy on the middle-signature locus]
\label{lem:routes-middle-weyl}
In the induced metric of a smooth Euclidean hypersurface, take principal
curvatures \((-a,-b,c,d)\), all four parameters positive. The three pair
sums are \(ab+cd,-ac-bd,-ad-bc\). For either Weyl endomorphism the squared
Hilbert--Schmidt norm is one twelfth of the sum of the three squared
differences between these sums. It is at least \((8/3)abcd\), with
equality precisely when \(a=b=c=d\). The scalar algebra and the Gauss
equation identifying the actual curvature are separate proof targets.
\end{lemma}

\begin{lemma}[Shape-operator deficit]
\label{lem:routes-shape-deficit}
For arbitrary real principal curvatures in the same block normalization,
\[
  \tfrac16 |A^\circ|^4-|W^\pm|_{\rm HS}^2
     \ \ge\ \tfrac89(-\det A)_+ .
\]
The tuple \((-1,3,3,3)\) realizes equality with negative determinant.
\end{lemma}

\begin{lemma}[Indexed Gauss mass]
\label{lem:routes-indexed-mass}
Let \(m_2(u)\) count the index-two critical points of the height in a
regular direction \(u\) for one smooth immersion of a closed oriented
homotopy four-sphere in \(\mathbb R^5\). With normalized spherical
probability, \(\omega=|S^4|\) and \(U\) the middle-signature locus,
\[
  \mathbb E[m_2]=\frac{2}{\omega}\int_U\det A\,dV_g .
\]
Both normal directions contribute. The area formula counts with
multiplicity on the smooth domain. Identify a measurable count and the
null exceptional directions before integrating.
\end{lemma}

\begin{proposition}[The topological floor for the middle count]
\label{prop:routes-middle-count}
On a nonstandard smooth homotopy four-sphere every Morse function has at
least two index-two critical points. The proposed formal argument uses
the actual extremal handlebody reductions, preserves the framed
attachments, and applies the one-middle-pair recognition theorem and
final disk gluing when the reduced count is at most one.
\end{proposition}

\begin{lemma}[Selecting a two-middle height]
\label{lem:routes-height-selection}
\lean{SmoothPoincare.prob_twoMiddle_gt_threeQuarters, SmoothPoincare.twoMiddle_probability_boundary}
\leanok
On a probability space let \(m\) be an integrable, measurable
natural-valued count. If \(m\ge2\) almost everywhere and
\(\mathbb E[m]<9/4\), then \(\Pr(m=2)>3/4\).
\end{lemma}

\begin{proof}
\leanok
Outside the two-count event the integer count is at least three almost
everywhere. Integrate this extra unit. A probability of one quarter
at three and three quarters at two gives equality at the excluded mean.
\end{proof}

\begin{corollary}[Heights of a subcritical immersion]
\label{cor:routes-subcritical-heights}
\uses{lem:routes-middle-weyl,lem:routes-indexed-mass,prop:routes-middle-count,lem:routes-height-selection}
If the induced class of the specified immersion on a nonstandard
candidate has \(E_+(g)<8\pi^2\), more than three quarters of directions
have exactly two middle critical points. The direction count uses no
Yamabe sign. Recognition still requires control of a selected diagram.
\end{corollary}

\begin{proposition}[A marked mixed cancellation]
\label{prop:routes-mixed-cancellation}
\uses{prop:routes-middle-count}
Suppose a reduced smooth handle presentation of \(M\) has one 0-handle,
one 1-handle, two 2-handles, one 3-handle and one 4-handle. Supply a
framed isotopy or a sequence of justified slides that preserves the
marked presentation and gives a transverse single geometric
intersection of one 2-handle attaching circle with the 1-handle belt
sphere. Smooth handle cancellation leaves at most one middle handle,
so \(M\) is standard. Producing this geometric intersection from a
selected height is open.
\end{proposition}

\begin{openproblem}[Coverage on each candidate smooth type]
\label{problem:routes-coverage}
\uses{prop:master-recognition,prop:master-immersion,cor:routes-subcritical-heights,prop:routes-mixed-cancellation,prop:master-short-cut,cor:master-repair-residual,prop:routes-joint-descent}
For each candidate, produce a recognized class, an immersed subcritical
height with a recognized marked diagram, or a positive-class minimizing
sequence on which a strict contradiction is realized. Quantitative
construction and coverage must refer to the same metric, immersion,
or chosen physical sequence in the implication that uses it.
\end{openproblem}

\subsection{Conditional comparisons on supplied witnesses}

\begin{proposition}[Measured replacements and two full cut sides]
\label{prop:ed03-measured-comparisons}
\uses{lem:master-competitor,prop:cut-marked-reassembly}
\lean{SmoothPoincare.geometricRepairedMass_le_excess_sum_of_pullback, SmoothPoincare.geometricRepairedMass_le_excess_sum, SmoothPoincare.geometricDiscardedMass_le_excess, SmoothPoincare.MarkedSideReplacement.reassemble}
\leanok
Fix the geometric positive-energy set of the original oriented smooth
four-manifold and use its physical chiral density in a specified frame.
For individually marked admissible repairs of known global cost,
\(\sum_j m_j\le N(e-I)+\sum_j\gamma_j\).
For two full measured cut sides, suppose their union is the source,
their intersection is volume-null, both caps are positive, and recognition
and smooth collar reassembly data are supplied for either choice of the
discarded cap. If \(e+\gamma_1+\gamma_2<8\pi^2\), a surviving side
\(s\) gives an oriented competitor with
\(m_{3-s}\le e-I+\gamma_s\). The resulting map carries the entire
opposite raw side onto the marked cap.
\end{proposition}
\begin{proof}
\leanok
Integrate the source curvature on the specified sets. For independent repairs,
pull each metric back and sum its comparison with the same infimum.
For the cut, the strict total bound selects a cap below the recognition
threshold. Use its corrected marking on a neighborhood to identify the
other cap with the source. Apply the infimum bound to that pullback.
\end{proof}

\begin{lemma}[A common Yamabe floor along an actual metric family]
\label{lem:ed03-path}
\uses{def:master-admissible}
\lean{SmoothPoincare.positiveAlongSubcriticalPath, SmoothPoincare.admissibleAlongSubcriticalMetricPath}
\leanok
Suppose a connected parameter family has a continuous sign detector,
positive initially, with the stated positive and zero sign identifications
for the class Yamabe constants of its metrics on the original type.
Supply integrable chiral densities and the normalized inequality
\(Y_t^2\ge48(8\pi^2-E_t)\).
If \(E_t\le8\pi^2-\delta\) throughout for one \(\delta>0\), every energy
is admissible and every \(Y_t\ge\sqrt{48\delta}\).
\end{lemma}
\begin{proof}
\leanok
At a zero of the sign detector the Euler inequality forces
\(E_t\ge8\pi^2\). The intermediate value theorem propagates the positive
sign. Take the square root of the Euler inequality and use each actual
framed metric to witness membership.
\end{proof}
